 
 \section{Mechanisms}
 \label{sec: mechanisms}

 
 \subsection{Potential Drivers of Long-Term Treatment Effects on Enrollment}\label{sec: mechanisms enrollment}




 
\subsubsection{Selective college match}\label{sec : college match}





PACE may have led students to enroll in more demanding programs, causing larger dropout among students from treated schools. To examine this channel, we begin by analyzing its effects on enrollment in STEM versus non-STEM programs at selective colleges, as STEM majors are typically more academically challenging. We assign value zero to both STEM and non-STEM enrollment for students who do not enroll in any selective college. Tables \ref{tab: TEenrolledovertimeSTEM} and \ref{tab: TEenrolledovertimetop15STEM} show that PACE increased enrollment rates in STEM and non-STEM fields almost equally, maintaining the same relative proportion between these fields as observed in the control group. The difference between treatment effects on STEM and non-STEM enrollment is always small and statistically insignificant across all years and subsamples, including the top 15\% of students at baseline where enrollment impacts are concentrated. Thus, PACE did not change students' relative propensity to pursue STEM versus non-STEM programs at selective colleges.\par 


Next, we examine the selectivity of degree programs chosen by students who enroll in selective colleges, where a degree program is defined as a college-major pair. We measure program selectivity using the average entrance exam score of regular-admission students, and set the outcome variable to missing for students who do not enroll in selective colleges. Columns (1) and (3) of Table \ref{tab: RCTselectivitydistanceselective}, which control for student characteristics, show no statistically significant differences in program selectivity between college entrants from treated and control schools.\footnote{ Lee bounds \citep{lee2009training}---which bound intensive margin impacts among always-enrollers---are large due to the substantial extensive margin effect on selective college enrollment, though they always include zero. We report them in Appendix Table \ref{tab: leeboundsselectivitydistanceonlyselective}.}\par 




We also examine the geographic distribution of enrollment by measuring the distance between students' high schools and their chosen selective college programs. Columns (2) and (4) of Table \ref{tab: RCTselectivitydistanceselective} show that, among selective college enrollees, students from treated schools enroll in programs that are closer to their high schools, though the difference is statistically insignificant. \par 



 The enrollment results align with the admission patterns, described in Appendix \ref{application_lists}. We do not find significant differences across treatment groups in the characteristics of the selective college programs to which students are admitted through the regular process (field of study, selectivity, and location). Similarly, there are no significant differences in admission patterns between the regular and PACE channels for students in PACE schools.\par 

 Finally, we examine whether college entrants from treated schools are more negatively selected on baseline measures of ability. Focusing on the sample of students in the top 15\% of their school at baseline---where college impacts are concentrated---we find that entrants from treated schools have lower grade 10 standardized test scores (-0.05 standard deviations; Table \ref{tabcollegeentrants}), but this difference is statistically insignificant. \par 

 These findings suggest that a worsening of the college match in terms of selectivity, field of study, or increased distance to college (and associated factors like separation from support networks or higher costs) is unlikely to explain the higher dropout rates among students from treated schools. Moreover, we find no statistically significant evidence of more negative selection on observed baseline ability, although we cannot rule out a more negative selection on unmeasured ability.





\subsubsection{Reductions in college preparedness}\label{sec:reductions_college_preparedness}
By reducing pre-college effort and achievement, PACE could have reduced college preparedness, leading to lower college persistence rates in the treatment group. To investigate this channel, we examine whether PACE lowered pre-college outcomes that predict persistence in college.\par 
 Appendix Table \ref{tab:persistenceSUA}  shows that, after controlling for student characteristics, GPA in the last high school year strongly predicts continuous enrollment or graduation five years after entering a selective college, independently of the entrance exam score and of the baseline test score (column (1)). GPA in the subjects tested on the entrance exam correlates more strongly with persistence than GPA in untested subjects (column (2)). The PSU score independently predicts persistence, although not significantly (columns (1) and (2)). If GPA and the PSU score at the end of high school are produced by a combination of baseline ability and study effort during high school, the administrative measure of baseline ability and our survey measure of study effort should both predict persistence. This is indeed what we find: both measures are significantly predictive, even after conditioning on the rich vector of student characteristics (columns (3) and (4)). \par 
 This evidence suggests that PACE reduced pre-college outcomes that predict persistence. Competence in the core high school subjects, which are tested on the entrance exam, seems to matter most for persistence. \par 


 \subsection{Potential Drivers of Negative Impacts on Pre-College Effort and Achievement}


 
 \subsubsection{Students' response to incentives}\label{sec:beliefs}


Preferential admissions introduce new admission requirements based on pre-college achievement. Since achievement is not a fixed trait but rather an outcome that responds to study effort, the introduction of new requirements can induce an endogenous response in study effort if students value college admission. Did students respond to incentives?\par 



\paragraph{Heterogeneity of impacts by absolute and relative ability.} To better understand the effort response, we examine effect heterogeneity along baseline within-school rank and baseline ability. We split the sample into quintiles of baseline ability and baseline within-school rank, and estimate the regression from equation \eqref{eq:main_OLS} on each sub-sample. The results are reported in Figure \ref{fig:hetero_all}. We do not find evidence of encouragement effects on pre-college effort or achievement, anywhere along the baseline relative and absolute ability distributions, and we find the negative impacts are spread across baseline relative and absolute ability.\par



\begin{figure}[ht!]
	\centering
	\includegraphics[width=0.85\linewidth]{Revision/Figures R/heterog_effects_effort_score.png}
	\vspace{-0.5cm}
	\caption{\label{fig:hetero_all} Heterogeneity of policy effects on pre-college effort and achievement. Notes: Each dot is the coefficient on {\itshape Treatment} from an OLS regression where: {\itshape Treatment} is a dummy variable indicating whether a student is in a school that was randomly assigned to be in the PACE program, the controls are the standard set of controls (see Figure \ref{fig:TEovertime}), Inverse Probability Weights and field-worker fixed effects are used, the estimation samples are quintiles in the within-school rank based on $9^{th}$ and $10^{th}$ grade GPA (left panel) and quintiles in the distribution of $10^{th}$ grade standardized test scores (right panel). The units of measurement of the treatment effects are standard deviations. The bars are $95\%$ confidence intervals built using standard errors clustered at the school level. Q-values for the family of quintiles in the upper left panel are $q(Q1)= 0.001$, $q(Q2)= 0.161$, $q(Q3)= 0.329$, $q(Q4)= 0.700$, $q(Q5)= 0.415$. The respective q-values for the upper right panel are $0.087, 0.412, 0.213, 0.412, 0.087$. The respective q-values for the lower left panel are $0.640, 0.875, 0.952, 0.105, 1.000$. The respective q-values for the lower right panel are $0.172, 0.172, 0.225, 0.345, 0.048$.}
	%\end{minipage}
\end{figure}
 
 
 
 

These patterns are hard to rationalize as a response to incentives under rational expectations. As shown theoretically in \citealp{bodoh2018college}, when students rationally respond to the incentives of percent rules, negative impacts are concentrated among those around the regular admission cutoff but well above the preferential admission cutoff. For these students the policy lowered returns to effort by guaranteeing an admission that was previously only within reach under sustained effort. Conversely, we would expect positive impacts among students near the top $15\%$ cutoff and for whom PACE brought within reach an admission that was previously unattainable. But these are not the patterns we find. \par  




A potential reason for not finding effects expected under rational expectations is that beliefs about own absolute and relative ability are systematically biased. Therefore, we examine students' beliefs next.\par 


 



\vspace{-10pt}


\paragraph{Students' misperceptions about their absolute and relative ability.} We elicited subjective expectations over the PSU entrance exam score using the survey question reported in the first row of Table \ref{tab:beliefelicitation}. The answers were given as ranges of the score. Using the midpoint of the range as the measure of perceived score, Table \ref{tab:summarybeliefs} shows that students display large over-optimism over their PSU entrance exam score (first two lines), on average expecting a score that is 0.6 standard deviations above the score they obtain. Figure \ref{fig: PSU_belief_distribution} confirms the large over-optimism by plotting the histograms of the raw survey answers and of the actual scores. 


 We elicited subjective expectations over own GPA and the top 15\% cutoff in the school using the survey questions reported in the second and third rows of Table \ref{tab:beliefelicitation}. Students display large over-optimism about their within-school rank, with over $40\%$ believing that their GPA is in the top $15\%$. While students hold accurate beliefs about their own GPA (GPA is measured on a scale from 1 to 7 and on average the GPA students expect differs from the one they obtain by less than 0.1 GPA points), they have a belief bias about the $85^{th}$ GPA percentile in their school of less than half GPA point (fourth row of the Table). This small belief bias in absolute terms is large in relative terms because of strong grade compression, that we document in Figures \ref{fig:GPA_hist} and \ref{fig:discrimination}.\footnote{First, we show that while grades can range from 1 to 7, the vast majority lie between 5 and 6.5. Second, we link grade data to baseline and end-line standardized achievement measures, and show that grades do not discriminate substantially among students of different baseline abilities, and much less than the end-line standardized achievement test we administered does.} These belief biases are consistent with the limited college experience of students' parents (over 90\% did not study beyond secondary education) and the lack of relative rank feedback in PACE schools.\par 

\par 
\input{Revision/Tables R/summary_beliefs}
\par 

Examining belief heterogeneity, Figure \ref{fig:hetero_beliefs} shows that students of all (absolute and relative) ability levels are over-optimistic; Table \ref{tab:correlatesbelieferrors} shows that belief biases are only weakly related to socioeconomic background in our homogeneously disadvantaged sample, with small correlations with parental education. The findings align with existing evidence that over-optimism is widespread in many contexts, including education (\citealp{stinebrickner2014major,hakimov2025confidence}).\par 
\paragraph{Response to perceived incentives.} As students have biased beliefs about their relative rank in the school and performance on the entrance exam, a natural question is whether the impacts on pre-college outcomes are consistent with a response to perceived rather than actual incentives.\par 

The belief patterns shown in Table \ref{tab:summarybeliefs} and Figure \ref{fig: PSU_belief_distribution} point to this mechanism. Students tend to expect their entrance exam scores to be near the national average, which is close to the regular admission cutoffs. Believing that admission is within reach, most students in the control sample prepare for the entrance exam (Table \ref{tab:summaryoutcomesbeliefs}). At the same time, students tend to perceive themselves as having a high within-school rank, which may lead them to believe that a preferential admission is guaranteed. On average, students view themselves as the type for whom PACE reduces the incentive to exert effort: those who are marginal for regular admission but confident of gaining preferential admission.

To further explore this channel, we examine the heterogeneity of impacts on pre-college outcomes by perceived absolute and relative ability. If students respond to perceived incentives, the marginal utility of effort is highest at the perceived top 15\% cutoff, as a small change in GPA can determine whether a student is in or out of the top 15\% and thus eligible for a preferential admission. The incentive to exert effort decreases as students perceive themselves to be further from this cutoff, leading to smaller treatment effects, even negative for those who perceive themselves to be well above the cutoff. Moreover, the negative impacts on effort should be strongest among students who are not confident of gaining regular admission, and believe they are above the top 15\% cutoff. \par


\input{Revision/Tables R/TE_pre_college_outcomes_perceived_dist_cutoff_f}

This is what we find, as shown in Table \ref{TEprecollegeoutcomesperceiveddistcutoff}. We define the perceived distance from the cutoff as the absolute value of the difference between a student's perceived own GPA and the perceived GPA of the $85^{th} $ percentile in their school. 
We regress pre-college outcomes on the treatment indicator, the controls, the measure of perceived distance from the cutoff, and its interaction with both the treatment indicator and controls. 
We perform these regressions on all students (Panel A) and on the sub-sample of students who believe they are in the top 15\% and at or below the median PSU score (Panel B).\footnote{The minimum PSU required for regular admission varies by program, averaging 480. The median subjective expectation for PSU scores falls within the 450-600 range. In Panel B, we exclude students above this median (i.e., those expecting a PSU between 600 and 850), as these students likely perceive themselves to be in a relatively secure position for regular admission.} As expected if students were responding to perceived incentives, the treatment effects on achievement, study effort, and GPA (both tested and untested subjects) diminish as students perceive themselves to be further from the cutoff. The coefficient on the interaction between treatment and perceived distance is negative for all pre-college outcomes and statistically significant for study effort and GPA in tested subjects. These effects are stronger for students who believe they are above the top 15\% cutoff but at or below the median PSU score (Panel B). For this subgroup, treatment effects are positive at the perceived cutoff (the treatment coefficient is positive but imprecisely estimated, reflecting the low density of students who perceive themselves to be at the cutoff) but become negative as the perceived GPA rises further above the perceived cutoff. The evidence, therefore, is consistent with a response to perceived incentives.\footnote{A caveat of these results is that subjective expectations were not elicited at the experiment's baseline, raising the concern that the expectations themselves could have been influenced by the treatment. In Appendix \ref{sec:robustnesshet} we provide robustness checks demonstrating this is unlikely to drive the results in Table \ref{TEprecollegeoutcomesperceiveddistcutoff}.}


Appendix \ref{sec:predvalbeliefs} shows that the belief measures have good predictive validity properties, giving us confidence in the findings presented in this section.



\subsubsection{Perceived college graduation likelihood and pre-college study effort}\label{sec:pgrad}

The evidence so far suggests that students on average lower their study effort because they perceive it is no longer needed to obtain a selective college admission. However, they would not do so if they believed pre-college effort was important to do well in college.\par 


We elicited students' beliefs about their likelihood of graduating from a selective college, if they were to enroll in one (Table \ref{tab:beliefelicitation}). We find that half of the students are certain they will graduate if admitted, and three quarters believe they have more than $50\%$ chance of graduating. Figure \ref{fig:p_grad_effect} shows that, despite its large impacts on pre-college effort, PACE had only limited impacts on this subjective belief, which was elicited after the effort reductions had occurred. Only 3.7 percent of the sample appear to be affected, not answering ``probably yes" when treated, opting instead for ``equally likely" (2 percent), ``probably not" (0.6 percent), and ``definitely not" (1.1 percent).\footnote{As is always the case with experimental data, we do not observe individual-level treatment effects, only averages. An alternative explanation to these patterns could be, for example, that 3.7 percent of the sample do not answer ``probably yes", opting instead for ``definitely yes", and an equal, offsetting fraction of the sample do not answer ``definitely yes", opting instead for ``equally likely" (2 percent), ``probably not" (0.6 percent), and ``definitely not" (1.1 percent). Any pattern that matches the net effects is consistent with the data. But these alternative, more convoluted explanations appear less likely.} Assigning numerical values to the survey answers reveals null effects on the average perceived graduation likelihood, irrespective of the regression specification (Appendix Table \ref{tab:pgradreg}). The heterogeneity analysis reported in Appendix Figure \ref{fig:heter_p_grad_effect} further shows that there was no substantial effect on the perceived graduation likelihood across baseline ability and within-school rank: regardless of the scale used to assign numerical values to the survey answers, the impacts hover around zero for most sub-samples, even among those who experienced considerable reductions in effort, achievement, or both, as per Figure \ref{fig:hetero_all}. Together, these empirical results suggest that students do not perceive pre-college effort as important for persistence in college.


\begin{figure}[ht!]
	\centering
	\includegraphics[width=0.65\linewidth]{Revision/Figures R/histogram_p_graduate_T_C.png}
	\vspace{-0.5cm}
	\caption{\label{fig:p_grad_effect} Distribution of survey responses regarding beliefs about the likelihood of graduating from a selective college conditional on enrolling, by treatment status (i.e., being in a PACE or control school). The figure includes $95\%$ confidence intervals for the difference between the proportion of treated and of control students giving each answer. The confidence intervals were obtained from standard errors clustered at high school level. An English translation of the survey question can be found in Appendix Table \ref{tab:beliefelicitation}.}
	%\end{minipage}
\end{figure}


\subsubsection{Other mechanisms: Teachers, schools, and perceived returns to college}\label{sec:other_mech}

Teachers can influence who obtains a preferential seat by adjusting their grading. If, in response to the percent plan policy, teachers manipulate their grading in ways that weaken the link between academic achievement and GPA, students in treated schools may have weaker incentives to study. This could help explain the observed reductions in pre-college effort. Teachers may also respond to the policy by changing their own effort or shifting the focus of instruction, which could affect student achievement both directly and indirectly—if, for instance, changes in teacher behavior influence how much students study. Schools might also adjust their academic support offerings, particularly for entrance exam preparation. However, evidence from supplementary teacher and principal surveys, along with merged data on grades and standardized test scores, suggests that these channels are unlikely to drive the decline in pre-college effort (Appendix \ref{sec:grading}).\par 

    If the light-touch orientation classes offered in PACE schools negatively affected students' beliefs about the net returns to college, they could have generated the reduction in pre-college study effort. \citealp{hastings2015effects} show that providing information on graduate earnings can change college applicants' choices in Chile. Although the orientation classes were not designed to provide information about returns, this remains an important channel to consider.\footnote{Perceived returns to college is an outcome that was not pre-specified in the pre-analysis plan. We added this outcome post hoc after observing declines in pre-college effort.} Evidence from survey data on perceived returns to college indicate that this channel is unlikely to drive the decline in pre-college effort: we find no effects on expected earnings nor on awareness of financial aid (Appendix \ref{sec:perceived_returns}).\par 




